\section{Conclusions}
\label{sec:conclusions}

We have measured the metallicity dependence of the main-sequence
mass--luminosity relation in the \gaia\ $G$ band from the relative motions of
14\,876 wide binaries within 1~kpc.  A modular Bayesian model infers the
shared metallicity of each binary from its two XP metallicities and the
PARSEC colour--magnitude relation, and propagates this posterior into a
hard-monotone, PARSEC-relative spline MLR constrained by a Rice-convolved
likelihood of the mass-scaled velocity statistic.  Our main conclusions are
the following.

\begin{enumerate}
\item With the velocity-distribution shape fixed to its orbital-population
reference, the inferred masses exceed PARSEC by $\approx30$\% near
$M_G\approx8$ at all metallicities.  The absolute dynamical mass scale is
controlled by the assumed intrinsic velocity distribution.
\item When the shape is sampled jointly with the MLR in four metallicity
bins, the MLR agrees with PARSEC to within $\approx15$\% for $M_G\le11$ in all
bins and to within $\approx7$\% for $M_G\le9$ in the near-solar bins, which
contain 96\% of the sample.
\item The MLR depends significantly on metallicity.  Relative to PARSEC,
stars with $\mh=0.3$--$0.6$ are $8$--$14$\% more massive over
$5\le M_G\le13$, and stars with $\mh=-1.0$--$-0.5$ are $3$--$6$\% less
massive for $M_G\le11$.
\item The metallicity sensitivity $\dd\ln M/\dd\mh$ rises from $0.27\pm0.04$
at $M_G=5$ to $0.75\pm0.09$ at $M_G=13$.  It has the sign predicted by
PARSEC but is $1.4$--$1.8$ times larger.
\item The preferred velocity distribution of the near-solar bins has a lower
cut-off and a broader core than the orbital-population reference and lies
at the edge of the explored parameter range.
\end{enumerate}

An absolute dynamical calibration of the MLR requires an independent
constraint on the intrinsic velocity distribution of wide binaries, for
example from forward modelling of the separation-dependent eccentricity
distribution and of the \gaia\ selection, together with a treatment of
hierarchical triples.  The same constraint would turn the metallicity
sensitivity measured here, which is currently conditional on the adopted
shape family, into a model-independent measurement, and the longer
astrometric baseline of \gaia\ DR4 will reduce the velocity uncertainties
that limit the metal-poor and metal-rich bins.
